\section{Reproducibility and Implementation Notes}
\label{app:reproducibility}
\setcounter{equation}{0}\setcounter{table}{0}\setcounter{figure}{0}
\subsection{Archive and reproducibility boundaries}
\label{app:evidence-map}
\WfourRepository
The evidence map is: relevance, \path{code/data/kappa_strength_summary.csv}; complete cell summaries, direct truth statistics and grid profiles, \path{code/results/main/n500/} and \path{code/results/main/n1000/}; false-candidate rejection, their \path{power_summary.csv} and compressed replication records; expanded-domain summaries and selected numerical comparisons, \path{code/results/diagnostics/}. The abstract's $0.942/3.900/0.094$ contrast uses direct truth acceptance, median grid-accepted measure and rejection at offset $+0.50$, respectively, for DML-IVQR at $n=1000$, $\tau=0.50$, $\kappa=0.10$; it is not a joint event or validated grid-set coverage. Historical file inventories are retained separately from the manuscript.
The selected package contains preserved scientific modules in \path{code/src/}, configuration and historical version records in \path{code/config/}, main outputs in \path{code/results/main/}, and diagnostic summaries in \path{code/results/diagnostics/}. The original environment is Windows R~3.4.3 with five workers, \texttt{quantreg}~5.34 and \texttt{hdm}~0.2.0; full version inventories and reporting instructions accompany the package rather than the manuscript. Main outputs remain frozen. Deterministic inspection of saved records differs from historical stochastic replay: worker-specific RNG streams and full main-run penalty states were not archived. Parent seed 675 alone does not establish replay or worker-stream independence. Large expanded-grid profiles, saved primitives and a portable historical runtime are not established as publicly deposited. The historical freeze identifier \texttt{ac9a9cca} remains unresolved and is not a verified release checkout.

Selected-case numerical audits found five Policy A/C decision differences among 24 saved candidate problems. Policy C retains the density/projection pathway but uses exact outcome residual ranks. These panels do not estimate an experiment-wide error rate or replace any historical results. Certificates concern fixed objectives and stated inputs, not every historical decision or adaptive calibration trajectory. Rights and both-author distribution approval remain unresolved; authenticated repository access does not establish anonymous public access or permission to redistribute.

\subsection{Historical evaluation sequence}
\label{app:evaluation-sequence}
One historical replication proceeds as follows.
\begin{enumerate}
\item Generate one set of primitives; construct a sample at each relevance setting. Within that setting all three estimators share $(Y,D,X,Z)$. Reusing primitives across settings does not hold realized $(Y,D)$ fixed.
\item For each quantile and estimator, profile the ordered 41-point primary grid. At each DML candidate $a$, call \texttt{lambda\_BC\_author} afresh (1000 artificial uniform columns), then fit \texttt{quantreg::rq(y-a*D\string~X, method="lasso")} with the resulting penalty and an intercept. Refit nuisances at each candidate; there is no post-Lasso refit. Oracle and Full instead call unpenalized \texttt{quantreg::rq} with their designated controls.
\item Form candidate outcome residuals and the proxy \texttt{dnorm(e, mean(e), var(e))}. Residualize instruments using the DML transformed regressions or benchmark matrix rule below; numerical termination and last-used penalty conventions are specified once in Appendix~\ref{app:finite-solvers}.
\item Apply the literal computed-residual rule $r_i=\mathbf1\{\widehat e_i\leq0\}$, with no residual tolerance. Set $s_i=\tau-r_i$, $g=\sum_i\widehat\Psi_i s_i$ and $S=\sum_i\widehat\Psi_i\widehat\Psi_i' s_i^2$. The statistic is $g'S^{-1}g$, equivalent to the normalized definition in Section~\ref{sec:framework}. Use direct \texttt{solve} and the retained nominal chi-square cutoff; singular or non-finite evaluations are failures, not acceptances.
\item Retain the first grid minimum for the point estimate and the stored grid decisions for informativeness summaries. After profiling, separately evaluate the true effect and six specified false candidates. DML again invokes its candidate routine, with fresh pivotal draws, in a separate five-worker evaluation. Direct-at-truth acceptance and false-candidate rejection therefore do not reuse stored grid statistics, even when candidates coincide. Failed profiles receive no point estimate; the documented failure conventions remain in force.
\end{enumerate}
The archived estimator includes both sample randomness and algorithmic penalty randomness. Common observed data do not imply common pivotal draws across estimators, candidates, or repeated calls. Worker-specific RNG stream assignments and the full main-run candidate penalty states were not archived as a replay specification; seed 675 alone neither reproduces every historical penalty nor establishes independence of worker streams. Saved Stage 2G cases with exported states have a narrower replay scope. The conventional Monte Carlo standard-error summaries do not verify historical stream independence. Frozen CSV inspection and selection of existing table/figure assets require no stochastic rerun; they do not recover missing RNG states.

\subsection{Penalty, density proxy, and residualization}
\label{app:penalty-implementation}
For each candidate, $b_0$ is the fitted outcome intercept and $b$ the slopes. In the preserved \texttt{rq.fit.lasso} augmented-observation implementation, the unscaled check-loss sum receives a penalty with factor $1/2$. The vector $\lambda^{rq}(a)$ is the supplied routine argument, including an intercept penalty. There is no post-Lasso refit. Let $s_k=(n^{-1}\sum_iX_{ik}^2)^{1/2}$; $U_{ir}$ are independent artificial Uniform$(0,1)$ draws, with $r=1,\ldots,1000$, and $\widehat Q_{0.9}$ is their empirical pivotal-statistic 90th percentile. Fresh draws are made on every penalty call. This is algorithmic randomness, not sample splitting.

In the following formulas, $\phi(e;\mu,\sigma)$ denotes a normal density parameterized by standard deviation. The proxy supplies $\widehat\sigma_\epsilon^2$ in that argument: the preserved R call is \texttt{dnorm(e, mean(e), var(e))}. It evaluates a fitted-normal density at residuals rather than estimating the theoretical conditional density at zero. The sample-variance convention is
\[
 \widehat\sigma_\epsilon^2(a)=\frac{1}{n-1}\sum_{i=1}^n
 \left[\widehat\epsilon_i(a)-\widehat\mu_\epsilon(a)\right]^2.
\]
This variance, not its square root, is supplied as the normal scale; the historical convention is described, not corrected. The instrument penalty is $\mathcal P_{j,a}(\gamma_j)=\sum_k\lambda^\Gamma_{jk}(a)|\gamma_{jk}|$, using default iterative heteroskedastic loadings. Its ordinary intercept is unpenalized; no post-Lasso refit is performed.
\begin{equation}
    \begin{pmatrix}
        \widehat b_0(a)\\
        \widehat\beta(a)
    \end{pmatrix}
    =
    \operatorname*{arg\,min}_{b_0,b}
    \left\{
        \sum_{i=1}^{n}
        \rho_{\tau}\!\left(
            Y_i-D_i^{\prime}a-b_0-X_i^{\prime}b
        \right)
        +
        \frac{1}{2}\left[
            \lambda_0^{rq}(a)|b_0|
            +\sum_{k=1}^{p}\lambda_k^{rq}(a)|b_k|
        \right]
    \right\}.
    \label{eq:beta-hat-dml}
\end{equation}

\begin{equation}
    \begin{aligned}
        T_r
        &=
        \max_{1\leq k\leq p}
        \left|
            \frac{
                \sum_{i=1}^{n}X_{ik}
                \left[\tau-\mathbf{1}\{U_{ir}<\tau\}\right]
            }{
                s_k\sqrt{\tau(1-\tau)}
            }
        \right|,\\
        \lambda^{rq}(a)
        &=
        2\,\widehat Q_{0.9}\!\left(\{T_r\}_{r=1}^{1000}\right)
        \sqrt{\tau(1-\tau)}
        (1,s_1,\ldots,s_p)^{\prime}.
    \end{aligned}
    \label{eq:bc-pivotal-penalty}
\end{equation}

\begin{equation}
    \begin{aligned}
        \widehat{\epsilon}_i(a)
        &=
        Y_i-D_i^{\prime}a-\widehat b_0(a)
        -X_i^{\prime}\widehat{\beta}(a),\\
        \widehat{\mu}_{\epsilon}(a)
        &=\frac{1}{n}\sum_{i=1}^{n}\widehat{\epsilon}_i(a),\\
        \widehat f_i(a)
        &=
        \phi\!\left(
            \widehat{\epsilon}_i(a);
            \widehat{\mu}_{\epsilon}(a),
            \widehat{\sigma}_{\epsilon}^{2}(a)
        \right),
    \end{aligned}
    \label{eq:estimated-profile-residual}
\end{equation}

\begin{equation}
    \begin{pmatrix}
        \widehat r_{j0}(a)\\
        \widehat\Gamma_{\mathrm{imp},j}(a)
    \end{pmatrix}
    =
    \operatorname*{arg\,min}_{(r_{j0},\gamma_j)}
    \left\{
        \sum_{i=1}^{n}
        \left[
            \sqrt{\widehat f_i(a)}Z_{ji}
            -r_{j0}
            -\sqrt{\widehat f_i(a)}X_i^{\prime}\gamma_j
        \right]^2
        +\mathcal{P}_{j,a}(\gamma_j)
    \right\}.
    \label{eq:dml-residualization-implementation}
\end{equation}

\begin{equation}
    \begin{aligned}
        \widehat\Psi_{ji}(a)
        &=
        Z_{ji}-\widehat r_{j0}(a)-X_i^{\prime}\widehat\Gamma_{\mathrm{imp},j}(a),\\
        \widehat\Psi_i(a)
        &=
        Z_i-\widehat r_0(a)-\widehat\Gamma_{\mathrm{imp}}(a)X_i.
    \end{aligned}
    \label{eq:estimated-residualized-instrument}
\end{equation}

\begin{equation}
    \begin{pmatrix}
        \widehat b_{0,\mathrm{orc}}(a)\\
        \widehat\beta_{\mathrm{orc}}(a)
    \end{pmatrix}
    =
    \operatorname*{arg\,min}_{b_0,b}
    \sum_{i=1}^{n}
    \rho_{\tau}\!\left(
        Y_i-D_i^{\prime}a-b_0-X_{S_0,i}^{\prime}b
    \right).
    \label{eq:oracle-beta}
\end{equation}

\begin{equation}
    \begin{pmatrix}
        \widehat b_{0,\mathrm{full}}(a)\\
        \widehat\beta_{\mathrm{full}}(a)
    \end{pmatrix}
    =
    \operatorname*{arg\,min}_{b_0,b}
    \sum_{i=1}^{n}
    \rho_{\tau}\!\left(
        Y_i-D_i^{\prime}a-b_0-X_i^{\prime}b
    \right).
    \label{eq:full-beta}
\end{equation}

\begin{equation}
    \begin{aligned}
        \widehat M_e(a)
        &=\frac{1}{n}\sum_{i=1}^{n}
          Z_iX_{e,i}^{\prime}\widehat f_{e,i}(a),&
        \widehat J_e(a)
        &=\frac{1}{n}\sum_{i=1}^{n}
          X_{e,i}X_{e,i}^{\prime}\widehat f_{e,i}(a),\\
        \widehat\delta_e(a)
        &=\widehat M_e(a)\widehat J_e(a)^{-1},&
        \widehat\Psi_{e,i}(a)
        &=Z_i-\widehat\delta_e(a)X_{e,i}.
    \end{aligned}
    \label{eq:benchmark-residualization}
\end{equation}
The first five equations describe DML-IVQR. The last three describe Oracle-GMM and Full-GMM, with $S_0=\{1,\ldots,10\}$, $X_{\mathrm{orc},i}=X_{S_0,i}$, and $X_{\mathrm{full},i}=X_i$. Their quantile regressions include an intercept, but their matrix residualization includes no constant regressor. DML retains the ordinary intercept of a transformed regression and subtracts it from the original instrument without dividing by $\sqrt{\widehat f_i(a)}$. It is therefore not the exact sample counterpart of the ideal density-weighted intercept. All procedures use same-sample nuisance fits and scores; cross-fitting is absent. The comparison does not isolate regularization from these residualization differences.

\subsection{Objectives and finite historical solver outputs}
\label{app:finite-solvers}
The argmin displays above specify optimization objectives; the archived coefficients are finite routine outputs, not certified exact minimizers. The DML outcome route is \texttt{quantreg::rq(..., method="lasso")} / \texttt{rq.fit.lasso} in \texttt{quantreg}~5.34. It uses augmented observations and the penalized intercept described above. Oracle and Full use the package's default unpenalized \texttt{rq} route. Computed residual signs, including exact floating-point zeros, feed the literal $\leq0$ rule; small objective errors alone do not establish rank or decision invariance.

Instrument fitting uses \texttt{hdm}~0.2.0 \texttt{rlasso}, \texttt{post=FALSE}, default heteroskedastic X-independent loadings and an unpenalized intercept. Outer loading calibration uses 15 iterations and tolerance $10^{-5}$; inner shooting uses a 1000-iteration budget, coefficient-change tolerance $10^{-5}$ and zero threshold $10^{-6}$. Early termination and thresholding produce finite outputs, not accuracy certificates. The reported next penalty need not equal the last-used penalty. Later refinement at the last-used penalty differs from recalibrating the adaptive trajectory; neither replaces the archived fits. Exact solver inventories and historical evidence paths remain in the preserved W3 implementation notes outside this manuscript.

\subsection{Retained density-divergence argument for the original DGP}
\label{app:density-divergence}
For the original $D=\Phi(T)$ construction, the retained argument is as follows. Put $\rho=0.3$, $q_\tau=\Phi^{-1}(\tau)$, $Q_j=Z_j-h_j(X)=z_j+v_j$, and $\mu_Q=\kappa(Q_1+Q_2)/2$. The observed instrument shifts $h_j(X)$ are given in the DGP. Since $Q\sim\mathcal N(0,2I_2)$, Gaussian conditioning in the retained derivation gives
\[
 T\mid Q\sim\mathcal N(\mu_Q,v_\kappa),\qquad v_\kappa=3-\kappa^2,
 \qquad
 T\mid u=q_\tau,Q\sim\mathcal N(\mu_Q+\rho q_\tau,s_\kappa^2),\quad
 s_\kappa^2=v_\kappa-\rho^2=2.91-\kappa^2.
\]
Controls are independent of $(T,u,Q)$, and conditioning on $(X,Z)$ identifies $Q$. For the structural residual $e_\tau=\Phi(T)(u-q_\tau)$, Bayes' formula and $u\perp(X,Z)$ yield the mixed-density identity
\begin{equation}
 f_{e_\tau\mid X,Z}(0)
 =\phi(q_\tau)\operatorname E\!\left[\frac{1}{\Phi(T)}\,\middle|\,u=q_\tau,Q\right].
 \label{eq:mixed-density-divergence}
\end{equation}
Here $s_\kappa^2>1$ at every investigated relevance setting (1.91 at $\kappa=1$ and 2.90 at $\kappa=0.10$). As $t\to-\infty$, $1/\Phi(t)\sim |t|/\phi(t)$. Against the conditional normal density, the integrand is proportional in its tail to
\[
 |t|\exp\!\left\{\frac{1-1/s_\kappa^2}{2}t^2+O(t)\right\},
\]
whose negative-tail integral diverges. Thus the original $f_{e_\tau\mid X,Z}(0)$ is infinite. Mixing conditional only on $X$ instead gives Gaussian variance $3-\rho^2=2.91>1$ and the same divergence. These are densities that integrate over treatment, not $f_{e_\tau\mid D,X,Z}(0)$ at a realized positive treatment value. Nor are they the historical fitted-normal proxy. This retained regularity limitation does not establish which feature causes any particular Monte Carlo rejection pattern and does not validate or repair nominal inference.

\subsection{Score algebra and its conditions}
\label{app:score-algebra}
The ideal nuisance vector and derivative identities are retained below. They presume finite density-weighted moments and permissible differentiation, conditions not established for the simulation DGP's non-finite true residual density. They are not implementation-level assertions. The executed proxy, residualization and nuisance fits do not inherit these properties automatically, and the chi-square cutoff is used nominally without an established validity result for the preserved implementations under this DGP. Same-sample fitting requires its own nuisance and generated-score justification; absence of cross-fitting alone does not prove invalidity.
\begin{equation}
    \eta(a)
    =
    \begin{pmatrix}
        c(a)\\
        \beta(a)\\
        \operatorname{vec}\{\delta(a)\}
    \end{pmatrix},
    \qquad
    \eta_0
    =
    \eta\!\left(\alpha_0(\tau)\right)
    =
    \begin{pmatrix}
        c_0(\tau)\\
        \beta_0(\tau)\\
        \operatorname{vec}\{\delta(\alpha_0(\tau))\}
    \end{pmatrix}.
    \label{eq:profile-eta}
\end{equation}

\begin{equation}
    \operatorname{E}\!\left[
        \left\{
            \tau-\mathbf{1}\!\left[
                Y_i\leq D_i^{\prime}a+c(a)+X_i^{\prime}\beta(a)
            \right]
        \right\}\widetilde X_i
    \right]
    =0.
    \label{eq:profile-beta-moment}
\end{equation}

\begin{equation}
    \begin{aligned}
        M_0
        &=\operatorname{E}\!\left[
            Z_i\widetilde X_i^{\prime}
            f_{\epsilon_0}(0\mid X_i,Z_i)
        \right],\\
        J_0
        &=\operatorname{E}\!\left[
            \widetilde X_i\widetilde X_i^{\prime}
            f_{\epsilon_0}(0\mid X_i,Z_i)
        \right],\\
        \delta_0&=M_0J_0^{-1}.
    \end{aligned}
    \label{eq:orthogonality-MJ}
\end{equation}

\begin{equation}
    \begin{aligned}
        \left.
        \partial_{(c,\beta^{\prime})}
        \operatorname{E}\!\left[
            g_{\tau}\!\left(W_i;\alpha_0(\tau),\eta\right)
        \right]
        \right|_{\eta=\eta_0}
        &=-\operatorname{E}\!\left[
            (Z_i-\delta_0\widetilde X_i)\widetilde X_i^{\prime}
            f_{\epsilon_0}(0\mid X_i,Z_i)
        \right]\\
        &=-\left[M_0-\delta_0J_0\right]
        =0.
    \end{aligned}
    \label{eq:orthogonality-beta}
\end{equation}

\begin{equation}
    \left.
    \partial_{\delta_j}
    \operatorname{E}\!\left[
        g_{\tau,j}\!\left(W_i;\alpha_0(\tau),\eta\right)
    \right]
    \right|_{\eta=\eta_0}
    =-\operatorname{E}\!\left[
        \left\{
            \tau-\mathbf{1}\!\left[
                Y_i\leq D_i^{\prime}\alpha_0(\tau)
                +c_0(\tau)
                +X_i^{\prime}\beta_0(\tau)
            \right]
        \right\}\widetilde X_i^{\prime}
    \right]
    =0.
    \label{eq:orthogonality-delta}
\end{equation}
The first derivative cancels by $\delta_0=M_0J_0^{-1}$; the second cancels by the profiled quantile moment. Orthogonality concerns local nuisance sensitivity, and does not establish identification or finite-sample calibration.
